% Luc Destombe --- licence CC BY-NC-SA 4.0
% https://creativecommons.org/licenses/by-nc-sa/4.0/deed.fr
% Fichier autonome : compiler avec pdflatex, deux fois (signets, nombre de pages).
\def\lyceeversion{1.0}
\documentclass[a4paper,10pt]{article}
\makeatletter
\newif\iflycee@niveaux
\lycee@niveauxtrue
\newif\iflycee@cm
\newif\iflycee@palatino

\RequirePackage[utf8]{inputenc}
\RequirePackage[T1]{fontenc}
\RequirePackage[french]{babel}
\RequirePackage{etoolbox}

\newcommand{\ShorthandsOff}{\@ifundefined{shorthandoff}{}{\shorthandoff{;:!?}}}
\newcommand{\ShorthandsOn}{\@ifundefined{shorthandon}{}{\shorthandon{;:!?}}}
\ShorthandsOff
\AtBeginDocument{\ShorthandsOn}
\AtBeginEnvironment{tikzpicture}{\ShorthandsOff}

\RequirePackage{geometry}
\RequirePackage{fancyhdr}
\RequirePackage{lastpage}
\RequirePackage{multicol}
\RequirePackage{array}
\RequirePackage{enumitem}
\RequirePackage{xcolor}
\RequirePackage{needspace}

\RequirePackage{amsmath,amssymb}
\iflycee@cm\else\RequirePackage{fourier}\fi
\DeclareMathAlphabet{\mathcal}{OMS}{cmsy}{m}{n}
\renewcommand{\ttdefault}{lmtt}
\begingroup\catcode`\_=\active \catcode`\^=\active
\global\let_\sb \global\let^\sp
\endgroup
\AtBeginDocument{\catcode`\_=12 \mathcode`\_="8000
  \catcode`\^=12 \mathcode`\^="8000 }
\newlength{\retraitcalculs}
\setlength{\retraitcalculs}{2em}

\RequirePackage{tikz}
\usetikzlibrary{arrows.meta,calc,babel,shapes.misc}
\RequirePackage[most]{tcolorbox}
\tcbuselibrary{skins,breakable}

\definecolor{coulDef}{HTML}{1F4E79}
\definecolor{coulProp}{HTML}{7030A0}
\definecolor{coulRem}{HTML}{C55A11}
\definecolor{coulEx}{HTML}{2E7D32}
\definecolor{coulExo}{HTML}{B03A2E}
\definecolor{coulPart}{HTML}{1F4E79}
\definecolor{coulMeth}{HTML}{1B6E4F}
\definecolor{coulCourbe}{HTML}{0B5ED7}
\definecolor{coulCourbeB}{HTML}{C00000}
\definecolor{coulCourbeC}{HTML}{2E7D32}
\definecolor{coulGrille}{HTML}{8C8C8C}
\definecolor{coulRep}{HTML}{1B6E4F}
\definecolor{coulBar}{HTML}{2E7D32}

\newcommand{\R}{\mathbb{R}}
\newcommand{\Rs}{\mathbb{R}^{*}}
\newcommand{\N}{\mathbb{N}}
\newcommand{\Ns}{\mathbb{N}^{*}}
\newcommand{\Z}{\mathbb{Z}}
\newcommand{\Q}{\mathbb{Q}}
\renewcommand{\le}{\leqslant}
\renewcommand{\ge}{\geqslant}
\renewcommand{\approx}{\simeq}
\newcommand{\vect}[1]{\overrightarrow{#1}}
\newcommand{\norme}[1]{\left\lVert #1 \right\rVert}
\newcommand{\intervalle}[2]{\left[#1\,;#2\right]}
\RequirePackage{cancel}
\renewcommand{\CancelColor}{\color{coulExo}}
\newcommand{\fois}[1]{\times{\color{coulExo}#1}}
\newcommand{\barre}[1]{\cancel{#1}}
\newlength{\lycee@eqh}\newlength{\lycee@eqd}
\newsavebox{\lycee@eqa}\newsavebox{\lycee@eqb}\newsavebox{\lycee@eqc}
\newcommand{\lycee@eqcol}[1]{\hbox{$\renewcommand{\arraystretch}{1.35}%
  \begin{array}[t]{@{}>{\displaystyle}r@{}>{\displaystyle{}}l@{}}#1\end{array}$}}
\newcommand{\lycee@eqmesure}[1]{%
  \ifdim\ht#1>\lycee@eqh\setlength{\lycee@eqh}{\ht#1}\fi
  \ifdim\dp#1>\lycee@eqd\setlength{\lycee@eqd}{\dp#1}\fi}
\newcommand{\lycee@eqtrait}{\hspace{1.2em}\tikz[baseline=0pt]\draw[dash pattern=on 2pt off 2pt,
  line width=0.5pt] (0,\lycee@eqh) -- (0,-\lycee@eqd);\hspace{1.2em}}
\newcommand{\eqproduit}[2]{\par\smallskip\noindent C'est une équation produit :
  $A\times B=0 \iff A=0$ ou $B=0$.\par\nopagebreak\smallskip
  \begingroup\setlength{\lycee@eqh}{0pt}\setlength{\lycee@eqd}{0pt}%
  \setbox\lycee@eqa=\lycee@eqcol{#1}\setbox\lycee@eqb=\lycee@eqcol{#2}%
  \lycee@eqmesure\lycee@eqa\lycee@eqmesure\lycee@eqb%
  \noindent\hspace*{\retraitcalculs}\usebox\lycee@eqa\lycee@eqtrait\usebox\lycee@eqb\endgroup\par\smallskip}
\newcommand{\eqproduittrois}[3]{\par\smallskip\noindent C'est une équation produit :
  $A\times B\times C=0 \iff A=0$ ou $B=0$ ou $C=0$.\par\nopagebreak\smallskip
  \begingroup\setlength{\lycee@eqh}{0pt}\setlength{\lycee@eqd}{0pt}%
  \setbox\lycee@eqa=\lycee@eqcol{#1}\setbox\lycee@eqb=\lycee@eqcol{#2}\setbox\lycee@eqc=\lycee@eqcol{#3}%
  \lycee@eqmesure\lycee@eqa\lycee@eqmesure\lycee@eqb\lycee@eqmesure\lycee@eqc%
  \noindent\hspace*{\retraitcalculs}\usebox\lycee@eqa\lycee@eqtrait\usebox\lycee@eqb\lycee@eqtrait\usebox\lycee@eqc\endgroup
  \par\smallskip}

\tikzset{grille/.style={coulGrille,line width=0.4pt}}
\newcommand{\quadrillage}[4]{\draw[grille,step=1] (#1,#2) grid (#3,#4);}
\tikzset{
  croix/.style={cross out,draw,line width=0.6pt,inner sep=0pt,minimum size=4pt},
  croixdroite/.style={croix,rotate=45,line width=0.8pt},
}
\newcommand{\pt}[3]{\path (#1) node[croix]{} node[#2,font=\small,inner sep=2.5pt]{$#3$};}

\newcommand{\lycee@repere}[4]{%
  \draw[grille,step=1] (#1,#2) grid (#3,#4);
  \draw[-{Stealth[length=2mm]},thick] (#1,0)--({#3+0.4},0) node[below left=-1pt]{$x$};
  \draw[-{Stealth[length=2mm]},thick] (0,#2)--(0,{#4+0.4}) node[below left=-1pt]{$y$};
  \node[below left,font=\scriptsize,inner sep=1.5pt] at (0,0) {$0$};}
\newcommand{\lycee@gradx}[1]{\foreach \k/\l in {#1}{%
  \draw (\k,0.07)--(\k,-0.07) node[below,font=\scriptsize,inner sep=1.5pt]{$\l$};}}
\newcommand{\lycee@grady}[1]{\foreach \k/\l in {#1}{%
  \draw (0.07,\k)--(-0.07,\k) node[left,font=\scriptsize,inner sep=1.5pt]{$\l$};}}
\newcommand{\lycee@intff}[2]{\left[#1\,;#2\right]}
\newcommand{\lycee@intfo}[2]{\left[#1\,;#2\right[}
\newcommand{\lycee@intof}[2]{\left]#1\,;#2\right]}
\newcommand{\lycee@intoo}[2]{\left]#1\,;#2\right[}
\newcommand{\lycee@ens}[1]{\left\{#1\right\}}
\AtBeginDocument{%
  \providecommand{\repere}{\lycee@repere}%
  \providecommand{\gradx}{\lycee@gradx}%
  \providecommand{\grady}{\lycee@grady}%
  \providecommand{\Cf}{\mathcal{C}\sb{f}}%
  \providecommand{\Cg}{\mathcal{C}\sb{g}}%
  \providecommand{\intff}{\lycee@intff}%
  \providecommand{\intfo}{\lycee@intfo}%
  \providecommand{\intof}{\lycee@intof}%
  \providecommand{\intoo}{\lycee@intoo}%
  \providecommand{\ens}{\lycee@ens}}

\newlist{questions}{enumerate}{3}
\setlist[questions,1]{label=\arabic*\textdegree),leftmargin=*,itemsep=2pt,topsep=3pt}
\setlist[questions,2]{label=\alph*),leftmargin=*,itemsep=1pt,topsep=2pt}
\setlist[questions,3]{label=\roman*),leftmargin=*,itemsep=1pt,topsep=1pt}
\setlist[itemize]{label=\textbullet,leftmargin=*,itemsep=2pt,topsep=3pt}

\newcommand{\besoinplace}[1]{\par\penalty\@M\begingroup
  \dimen@=#1\relax\dimen@ii=\pagegoal\advance\dimen@ii-\pagetotal
  \ifdim\dimen@>\dimen@ii\ifdim\dimen@ii>\z@\vfil\fi\break\fi\endgroup}

\newcommand{\lycee@pied}{\thepage/\pageref{LastPage}}
\newcommand{\entete}[2]{%
  \pagestyle{fancy}\fancyhf{}%
  \fancyhead[L]{\footnotesize\color{coulPart}\textbf{#1}}%
  \fancyhead[R]{\footnotesize\color{coulPart}\textbf{#2}}%
  \fancyfoot[C]{\footnotesize\lycee@pied}%
  \renewcommand{\headrulewidth}{0.4pt}%
  \renewcommand{\headrule}{\hbox to\headwidth{%
    \color{coulPart!40}\leaders\hrule height \headrulewidth\hfill}}}

\RequirePackage{ccicons}
\newcommand{\lycee@licence}{{\scriptsize\color{gray}%
  \copyright\ Luc Destombe --- \ccbyncsa\ CC BY-NC-SA 4.0}}
\newif\iflycee@licence \lycee@licencetrue
\newcommand{\sanslicence}{\lycee@licencefalse}
\AtBeginDocument{\iflycee@licence\AddToHookNext{shipout/foreground}{%
  \put(0,0){\rlap{\kern\dimexpr1in+\hoffset+\oddsidemargin+\textwidth\relax
    \llap{\raisebox{-\dimexpr1in+\voffset+\topmargin+\headheight+\headsep
      +\textheight+\footskip\relax}{\lycee@licence}}}}}\fi}

\newcommand{\lycee@site}{faire-des-maths.fr}
\newcommand{\lycee@version}{{\scriptsize\color{gray}\lycee@site
  \ifdefined\lyceeversion\ --- Version~\lyceeversion\fi}}
\AtBeginDocument{\iflycee@licence\AddToHookNext{shipout/foreground}{%
  \put(0,0){\rlap{\kern\dimexpr1in+\hoffset+\oddsidemargin\relax
    \raisebox{-\dimexpr1in+\voffset+\topmargin+\headheight+\headsep
      +\textheight+\footskip\relax}{\lycee@version}}}}\fi}

\newcommand{\formulesserrees}{%
  \setlength{\abovedisplayskip}{3pt}\setlength{\belowdisplayskip}{3pt}%
  \setlength{\abovedisplayshortskip}{2pt}\setlength{\belowdisplayshortskip}{3pt}}

\setlength{\parindent}{0pt}

\RequirePackage[implicit=false,hidelinks,pdfencoding=auto,psdextra,bookmarksopen,
  bookmarksopenlevel=1,pdfcreator={},pdfproducer={}]{hyperref}
\RequirePackage{bookmark}
\hypersetup{pdfauthor={Luc Destombe},pdfkeywords={CC BY-NC-SA 4.0}}
\newcounter{lycee@signet}
\newcommand{\signet}[3][]{\stepcounter{lycee@signet}%
  \edef\lycee@dest{\if\relax\detokenize{#1}\relax
    signet.\arabic{lycee@signet}\else#1\fi}%
  \hypertarget{\lycee@dest}{}\bookmark[level=#2,dest=\lycee@dest]{#3}}
\pdfstringdefDisableCommands{%
  \def\R{R}\def\N{N}\def\Z{Z}\def\Q{Q}\def\vect#1{#1}\def\Cf{Cf}\def\Cg{Cg}\def\fois#1{×#1}%
  \def\barre#1{#1}\def\intervalle#1#2{[#1;#2]}\def\intff#1#2{[#1;#2]}%
  \def\textsuperscript#1{#1}\def\dfrac#1#2{#1/#2}\def\frac#1#2{#1/#2}%
  \def\vec#1{#1⃗}}%
\RequirePackage[official]{eurosym}

\geometry{top=0.8cm,bottom=1.3cm,left=1cm,right=1cm,footskip=0.6cm}

\pagestyle{fancy}
\fancyhf{}
\renewcommand{\headrulewidth}{0pt}
\fancyfoot[C]{\footnotesize Page \thepage{} / \pageref{LastPage}}

\newcommand{\titrefiche}[2]{%
  \begin{center}
    \Large\textbf{#1}\\[2pt]
    \large\textbf{#2}\\[2pt]
  \end{center}
  \vspace{0.10cm}}

\setlength{\columnseprule}{0.4pt}
\setlength{\columnsep}{15pt}
\renewcommand{\columnseprulecolor}{\color{black!45}}
\setlength{\multicolsep}{2pt}

\newcounter{partiefiche}
\renewcommand{\thepartiefiche}{\Roman{partiefiche}}
\newif\ifapresbandeau
\newcommand{\lycee@bandeau}[1]{%
  \par\penalty-600
  \vspace{0.08cm}%
  \noindent\colorbox{coulPart}{%
    \parbox{\dimexpr\linewidth-2\fboxsep\relax}{%
      \color{white}\bfseries\raggedright\strut#1\strut}}%
  \nopagebreak\par\nopagebreak
  \vspace{0.06cm}\nopagebreak
  \apresbandeautrue}
\newcommand{\partiefiche}[1]{%
  \refstepcounter{partiefiche}%
  \lycee@bandeau{\signet{1}{\thepartiefiche{} --- #1}\thepartiefiche{} --- #1}}
\newcommand{\partiefichesansnum}[1]{\lycee@bandeau{\signet{1}{#1}#1}}

\newcounter{exo}
\newlength{\espaceexo}\setlength{\espaceexo}{7pt}
\newlength{\espacetitre}\setlength{\espacetitre}{2pt}
\newcommand{\exo}[1][\relax]{%
  \par
  \ifapresbandeau\apresbandeaufalse\else\penalty-150\addvspace{\espaceexo}\fi
  \refstepcounter{exo}%
  \noindent\signet[exercice-\arabic{exo}]{2}{Exercice \arabic{exo}}%
  \textbf{\color{coulExo}\underline{Exercice \theexo}}%
  \ifx\relax#1\relax\else\textbf{ : #1}\fi
  \par\nopagebreak\vspace{\espacetitre}\nopagebreak
  \ignorespaces}

\setlist[enumerate]{nosep,leftmargin=*,itemsep=2pt,topsep=2pt}
\setlist[itemize]{nosep,leftmargin=*,topsep=2pt}
\newcommand{\choix}[3]{\par\hspace*{1em}%
  \makebox[0.3\linewidth][l]{a) #1}\makebox[0.3\linewidth][l]{b) #2}c) #3}
\newcommand{\deux}[2]{\par\hspace*{0.6em}%
  \makebox[0.48\linewidth][l]{#1}#2}
\newcommand{\trois}[3]{\par\hspace*{0.6em}%
  \makebox[0.32\linewidth][l]{#1}\makebox[0.32\linewidth][l]{#2}#3}

  \renewcommand{\exo}[1][\relax]{%
    \ifapresbandeau\apresbandeaufalse\else\par\penalty-150\fi
    \refstepcounter{exo}%
    \vspace{0.05cm}%
    \noindent\signet[exercice-\arabic{exo}]{2}{Exercice \arabic{exo}}%
    \textbf{\color{coulExo}\underline{Exercice \theexo}}%
    \ifx\relax#1\relax\else\textbf{ : #1}\fi%
    \par\nopagebreak
    \ignorespaces}
  \newcommand{\rappel}[1]{%
    \par\vspace{0.06cm}\nopagebreak
    \noindent\fcolorbox{black!35}{black!5}{%
      \parbox{\dimexpr\linewidth-2\fboxsep-2\fboxrule\relax}{%
        \small\textbf{Rappel.} #1}}%
    \par\vspace{0.06cm}\nopagebreak}
  \setlist[enumerate]{nosep,leftmargin=*}
  \setlist[itemize]{nosep,leftmargin=*}
  \setlength{\multicolsep}{3pt plus 1pt minus 1pt}
  \setlength{\premulticols}{10pt}
  \setlength{\postmulticols}{10pt}
  \newlist{expr}{enumerate}{1}
  \setlist[expr]{label={\Alph*\ $=$},nosep,leftmargin=*,itemsep=0pt}

\renewenvironment{center}{\par\addvspace{3pt}\centering}{\par\addvspace{3pt}}
\formulesserrees
\AtBeginDocument{\formulesserrees}
\clubpenalty=10000
\widowpenalty=10000
\displaywidowpenalty=10000
\brokenpenalty=10000
\setlength{\parskip}{0pt}

\RequirePackage{ragged2e}
\setlength{\RaggedRightRightskip}{0pt plus 1fil}
\AtBeginDocument{\RaggedRight\binoppenalty=10000 \relpenalty=10000 }
\g@addto@macro\@parboxrestore{\RaggedRight}
\makeatother

\providecommand{\C}{\mathbb{C}}
\ShorthandsOff
\tikzset{
  vecteurB/.style={-{Stealth[length=2.4mm]},line width=1.1pt,coulCourbeB},
  unite/.style={-{Stealth[length=2mm]},line width=1.2pt,black},
}
\newcommand{\planc}[4]{%
  \quadrillage{#1}{#2}{#3}{#4}
  \draw[-{Stealth[length=2mm]},thick] (#1,0)--({#3+0.4},0);
  \draw[-{Stealth[length=2mm]},thick] (0,#2)--(0,{#4+0.4});
  \draw[unite] (0,0)-- node[midway,below,font=\footnotesize,inner sep=1.5pt]{$\vec{u}$} (1,0);
  \draw[unite] (0,0)-- node[midway,left,font=\footnotesize,inner sep=1.5pt]{$\vec{v}$} (0,1);
  \node[below left,font=\scriptsize,inner sep=1.5pt] at (0,0) {$O$};}
\ShorthandsOn

\begin{document}

\titrefiche{1\textsuperscript{re} STI2D --- Nombres complexes}{Fiche d'exercices du chapitre}

\begin{multicols}{2}\raggedcolumns

\partiefiche{L'ensemble des nombres complexes}

\exo[Partie réelle, partie imaginaire]
Pour chacun des nombres complexes ci-dessous, donner sa partie réelle et sa partie
imaginaire ; dire s'il est un nombre réel ou un imaginaire pur.
\begin{center}
$z_{1}=6+2i$ \quad $z_{2}=-4-i$ \quad $z_{3}=-7i$ \\
$z_{4}=3$ \quad $z_{5}=5i-1$ \quad $z_{6}=i$
\end{center}

\exo[En phrase complète]
Pour chacun des nombres $w_{1}=-3+8i$, $w_{2}=9i$ et $w_{3}=-11$, écrire une phrase
complète du type : « La partie réelle de $w_{1}$ est \dots{} et sa partie imaginaire
est \dots{} ».

\exo[Retrouver le nombre]
Écrire sous forme algébrique le nombre complexe $z$ tel que :
\begin{enumerate}[label=\alph*)]
  \item $\mathrm{Re}(z)=4$ et $\mathrm{Im}(z)=-3$ ;
  \item $\mathrm{Re}(z)=0$ et $\mathrm{Im}(z)=2$ ;
  \item $\mathrm{Re}(z)=-1$ et $\mathrm{Im}(z)=0$ ;
  \item $\mathrm{Re}(z)=\frac{1}{2}$ et $\mathrm{Im}(z)=-\frac{1}{3}$.
\end{enumerate}

\exo[Nombre réel ou imaginaire pur]
Soit $x$ un nombre réel, $z=(x-2)+5i$ et $z'=4+(2x+6)i$.
\begin{enumerate}[label=\alph*)]
  \item Pour quelle valeur de $x$ le nombre $z$ est-il un imaginaire pur ?
  \item Pour quelle valeur de $x$ le nombre $z'$ est-il un nombre réel ?
\end{enumerate}

\partiefiche{Somme et produit}

\exo[Sommes et produits]
Soit $z_{1}=4+i$ et $z_{2}=2-3i$. Écrire sous forme algébrique :
\begin{center}
{\renewcommand{\arraystretch}{1.5}
\begin{tabular}{@{}l@{\hspace{2.5em}}l@{\hspace{2.5em}}l@{}}
$A=z_{1}+z_{2}$ & $B=z_{1}-z_{2}$ & $C=z_{1}z_{2}$\\
$D=z_{1}^{2}$ & $E=2z_{1}-3z_{2}$ &
\end{tabular}}
\end{center}

\exo[Forme algébrique]
Écrire sous forme algébrique :
\begin{center}
{\renewcommand{\arraystretch}{1.5}
\begin{tabular}{@{}l@{\hspace{2.5em}}l@{}}
$A=(2-5i)+(-6+3i)$ & $B=7i-4-(1-2i)$\\
$C=(-2+i)(3+4i)$ & $D=(5-i)(5+i)$\\
$E=(1+3i)^{2}$ & $F=(2-i)^{2}$
\end{tabular}}
\end{center}

\exo[Puissances de $i$]
\begin{enumerate}[label=\alph*)]
  \item Sachant que $i^{3}=i^{2}\times i$, écrire $i^{3}$, $i^{4}$, $i^{5}$ et $i^{6}$
    sous forme algébrique.
  \item En déduire $(2i)^{3}$.
  \item Calculer $(1+i)^{2}$, puis $(1+i)^{4}$.
\end{enumerate}

\exo[Impédances en série]
En électricité, un dipôle se modélise par son impédance complexe $Z=R+iX$, en ohms :
$R$ est sa résistance et $X$ sa réactance. Deux dipôles en série ont pour impédance la
somme de leurs impédances. Soit $Z_{1}=30+40i$ et $Z_{2}=20-10i$.
\begin{enumerate}[label=\alph*)]
  \item Donner la résistance et la réactance du premier dipôle.
  \item Calculer l'impédance $Z_{1}+Z_{2}$ des deux dipôles en série.
  \item Un troisième dipôle a pour impédance $Z_{3}=10+Xi$. Pour quelle valeur de $X$
    l'impédance $Z_{1}+Z_{3}$ est-elle un nombre réel ?
\end{enumerate}

\exo[Vérifier une égalité]
Soit $z=1+2i$.
\begin{enumerate}[label=\alph*)]
  \item Calculer $z^{2}$.
  \item Vérifier que $z^{2}-2z+5=0$.
  \item Vérifier que le nombre $w=1-2i$ vérifie aussi $w^{2}-2w+5=0$.
\end{enumerate}

\partiefiche{Le plan complexe}

\exo[Lire et placer]
Dans le repère ci-dessous :
\begin{enumerate}[label=\alph*)]
  \item lire les affixes des points $A$, $B$, $C$, $D$, $E$ et du vecteur $\vec{w}$ ;
  \item placer les points $F(-3+3i)$, $G(4i)$, $H(-1)$ et $K(2-3i)$.
\end{enumerate}
\begin{center}
\begin{tikzpicture}[x=0.5cm,y=0.5cm]
  \planc{-4}{-4}{4}{4}
  \gradx{-3,-2,-1,2,3} \grady{-3,-2,-1,2,3}
  \pt{2,3}{above right}{A} \pt{-3,1}{above left}{B} \pt{3,0}{above right}{C}
  \pt{0,-2}{right}{D} \pt{-2,-3}{below left}{E}
  \draw[vecteurB] (1,-1)-- node[midway,right,font=\footnotesize]{$\vec{w}$} (3,-4);
\end{tikzpicture}
\end{center}

\exo[Sur les axes]
Dans un repère orthonormé direct, placer les points $A(1+2i)$, $B(-3+i)$, $C(4)$,
$D(-2i)$, $E(-1-3i)$ et $F(3i)$. Quels points ont pour affixe un nombre réel ? un
imaginaire pur ? Répondre par une phrase.

\exo[Affixe d'un vecteur]
Dans un repère orthonormé direct, placer le point $A(-1+i)$, puis tracer le
représentant d'origine $A$ du vecteur $\vec{w}$ d'affixe $3-2i$. Quelle est l'affixe
de son extrémité ?

\exo[Droites du plan complexe]
\begin{enumerate}[label=\alph*)]
  \item Placer trois points dont l'affixe a pour partie réelle $2$. Où sont tous ces
    points ?
  \item Même question avec une partie imaginaire égale à $-1$.
  \item Calculer $(2-i)^{2}$, puis placer le point $K$ d'affixe $(2-i)^{2}$.
\end{enumerate}

\partiefiche{Conjugué d'un nombre complexe}

\exo[Conjugués]
\begin{enumerate}[label=\alph*)]
  \item Donner le conjugué de chacun des nombres : $z_{1}=7+2i$, $z_{2}=3i-4$,
    $z_{3}=1-i$, $z_{4}=-9$, $z_{5}=-5i$, $z_{6}=i$.
  \item Placer le point $M$ d'affixe $z_{2}$ et le point $M'$ d'affixe
    $\overline{z_{2}}$.
\end{enumerate}

\exo[Produit $z\overline{z}$]
Pour chaque nombre $z_{1}=4-i$, $z_{2}=-3+5i$, $z_{3}=2i$ et $z_{4}=6$, calculer
$z\overline{z}$, puis écrire une phrase : « Le conjugué de \dots{} est \dots{} ».

\exo[Deux façons]
Soit $z_{1}=2+i$ et $z_{2}=1-3i$.
\begin{enumerate}[label=\alph*)]
  \item Calculer $z_{1}+z_{2}$, puis son conjugué.
  \item Calculer $\overline{z_{1}}+\overline{z_{2}}$. Que remarque-t-on ?
  \item Calculer $\overline{z_{1}z_{2}}$ de deux façons.
\end{enumerate}

\exo[Somme et différence avec le conjugué]
Soit $z_{1}=4-3i$ et $z_{2}=-1+2i$.
\begin{enumerate}[label=\alph*)]
  \item Pour chaque nombre $z_{1}$, $z_{2}$, calculer $z+\overline{z}$ et
    $z-\overline{z}$.
  \item Lesquels de ces résultats sont des nombres réels ? des imaginaires purs ?
  \item Placer les points d'affixes $z_{1}$ et $\overline{z_{1}}$.
\end{enumerate}

\partiefiche{Quotient de deux nombres complexes}

\exo[Quotients]
Écrire sous forme algébrique :
\begin{center}
$z_{1}=\dfrac{5}{1-2i}$ \quad $z_{2}=\dfrac{2+6i}{1+i}$ \quad
$z_{3}=\dfrac{1+5i}{i}$ \quad $z_{4}=\dfrac{1}{4+i}$
\end{center}

\exo[Quotients, suite des calculs]
Écrire sous forme algébrique :
\begin{center}
$w_{1}=\dfrac{4}{3i}$ \quad $w_{2}=\dfrac{5+2i}{3i}$ \quad
$w_{3}=\dfrac{1+i}{1-i}$ \quad $w_{4}=\dfrac{3+4i}{1+2i}$
\end{center}

\exo[Conjugué d'un quotient]
Soit $z=\dfrac{3-i}{2+i}$.
\begin{enumerate}[label=\alph*)]
  \item Écrire $z$ sous forme algébrique, puis donner $\overline{z}$.
  \item Écrire $\dfrac{3+i}{2-i}$ sous forme algébrique. Quelle propriété du cours
    retrouve-t-on ?
\end{enumerate}

\exo[Impédances en parallèle]
Une résistance d'impédance $Z_{1}=10$ et une bobine d'impédance $Z_{2}=10i$ sont
montées en parallèle. L'impédance du montage est $Z=\dfrac{Z_{1}Z_{2}}{Z_{1}+Z_{2}}$.
\begin{enumerate}[label=\alph*)]
  \item Calculer $Z_{1}Z_{2}$ et $Z_{1}+Z_{2}$.
  \item En déduire $Z$ sous forme algébrique.
  \item Donner la partie réelle et la partie imaginaire de $Z$.
\end{enumerate}

\partiefiche{Module d'un nombre complexe}

\exo[Modules]
Soit $z_{1}=6+8i$, $z_{2}=-3i$, $z_{3}=5-12i$, $z_{4}=-4$, $z_{5}=1+i$ et
$z_{6}=\sqrt{7}+3i$.
\begin{enumerate}[label=\alph*)]
  \item Calculer le module de chacun de ces nombres.
  \item En déduire $|z_{1}z_{3}|$ et $\left|\dfrac{z_{3}}{z_{1}}\right|$.
\end{enumerate}

\exo[Module d'un produit, d'un quotient]
Calculer, à l'aide des propriétés du module :
\begin{center}
$A=|(3-4i)(1+2i)|$ \quad $B=\left|\dfrac{1+i}{3-4i}\right|$ \\
$C=|2i(3+4i)|$ \quad $D=\left|(1+i)^{4}\right|$
\end{center}

\exo[Points d'un cercle]
Soit les points $A(3+4i)$, $B(-5)$, $C(4-3i)$ et $D(5i)$.
\begin{enumerate}[label=\alph*)]
  \item Calculer $OA$, $OB$, $OC$ et $OD$.
  \item Que peut-on dire de ces quatre points ? Répondre par une phrase.
  \item Le point $A'$ d'affixe $\overline{z_{A}}$ est-il sur le même cercle ?
\end{enumerate}

\exo[Intensité dans un dipôle]
Un dipôle a pour impédance $Z=30+40i$ (en ohms). Sous une tension de $230$~V,
l'intensité qui le traverse vaut $I=\dfrac{230}{|Z|}$ (en ampères).
\begin{enumerate}[label=\alph*)]
  \item Donner la partie réelle et la partie imaginaire de $Z$.
  \item Calculer $|Z|$.
  \item En déduire l'intensité $I$.
\end{enumerate}

\partiefiche{Affixe d'un vecteur et distance}

\exo[Nature d'un triangle]
Soit les points $A(1+i)$, $B(4+i)$ et $C(1+5i)$.
\begin{enumerate}[label=\alph*)]
  \item Calculer les affixes des vecteurs $\vect{AB}$, $\vect{AC}$ et $\vect{BC}$.
  \item En déduire les longueurs $AB$, $AC$ et $BC$.
  \item Le triangle $ABC$ est-il isocèle ? Est-il rectangle ?
\end{enumerate}

\exo[Centre d'un cercle]
Soit les points $M(1+i)$, $A(4+5i)$, $B(-4+i)$ et $C(1-4i)$.
\begin{enumerate}[label=\alph*)]
  \item Calculer $MA$, $MB$ et $MC$.
  \item En déduire que $A$, $B$ et $C$ sont sur un même cercle, dont on donnera le
    centre et le rayon.
\end{enumerate}

\exo[Trajet d'un drone]
Sur le plan d'un chantier, une unité vaut $100$~m. Un drone part du point $D(-3+2i)$
et vole en ligne droite jusqu'au point $E(5-4i)$.
\begin{enumerate}[label=\alph*)]
  \item Calculer l'affixe du vecteur $\vect{DE}$, puis la distance parcourue, en mètres.
  \item Le drone va ensuite au point $F$ d'affixe $\overline{z_{E}}$. Donner $z_{F}$
    et calculer la distance $EF$, en mètres.
  \item Il revient enfin en $D$. Calculer la longueur totale du trajet, arrondie au
    mètre.
\end{enumerate}

\partiefiche{Milieu et parallélogramme}

\exo[Parallélogramme]
Soit les points $A(1+i)$, $B(4+2i)$, $C(3+5i)$ et $D(4i)$.
\begin{enumerate}[label=\alph*)]
  \item Démontrer que $ABCD$ est un parallélogramme.
  \item Calculer l'affixe de son centre $K$.
\end{enumerate}

\exo[Milieux]
Calculer l'affixe du milieu de chacun des segments $[AB]$, $[CD]$ et $[EF]$, avec
$A(3+2i)$, $B(1-6i)$, $C(-5+i)$, $D(2+4i)$, $E(4i)$ et $F(-3)$.

\exo[Quatrième sommet]
Soit les points $A(-1+2i)$, $B(3+i)$ et $C(4+4i)$.
\begin{enumerate}[label=\alph*)]
  \item Écrire l'égalité des affixes qui traduit que $ABCD$ est un parallélogramme.
  \item En déduire l'affixe du point $D$.
  \item Calculer l'affixe du centre de $ABCD$.
  \item Placer les quatre points dans un repère.
\end{enumerate}

\exo[Losange]
Soit les points $A(3+i)$, $B(1+3i)$ et $C(4+4i)$.
\begin{enumerate}[label=\alph*)]
  \item Démontrer que $OACB$ est un parallélogramme.
  \item Calculer $OA$ et $OB$. En déduire la nature de $OACB$.
  \item Calculer l'affixe de son centre.
\end{enumerate}

\partiefiche{Argument d'un nombre complexe}

\exo[Lire module et argument]
Pour chacun des points de la figure, donner le module et un argument de son affixe,
puis l'écrire en phrase pour $E$ : « Le point $E$ est à \dots{} de $O$ et l'angle
$(\vec{u}\,;\vect{OE})$ mesure \dots{} ».
\begin{center}
\begin{tikzpicture}[x=0.5cm,y=0.5cm]
  \planc{-4}{-4}{4}{4}
  \gradx{-3,-2,-1,2,3} \grady{-3,-2,-1,2,3}
  \pt{3,0}{above right}{A} \pt{-2,0}{above left}{B} \pt{0,3}{right}{C}
  \pt{0,-1}{right}{D} \pt{3,3}{above right}{E} \pt{-2,-2}{below left}{F}
  \pt{1,-1}{below right}{G} \pt{-3,3}{above left}{H}
\end{tikzpicture}
\end{center}

\exo[Placer avec module et argument]
Dans un repère orthonormé direct d'unité $1$~cm, placer à la règle, au compas et au
rapporteur les points :
\begin{center}
$L$ : $|z_{L}|=2$ et $\arg(z_{L})=\frac{\pi}{3}$ \quad
$M$ : $|z_{M}|=3$ et $\arg(z_{M})=-\frac{\pi}{6}$ \\
$N$ : $|z_{N}|=4$ et $\arg(z_{N})=\frac{\pi}{2}$ \quad
$R$ : $|z_{R}|=3$ et $\arg(z_{R})=\pi$ \\
$S$ : $|z_{S}|=2$ et $\arg(z_{S})=\frac{3\pi}{4}$ \quad
$T$ : $|z_{T}|=1$ et $\arg(z_{T})=-\frac{2\pi}{3}$
\end{center}

\exo[Radar]
Un radar, placé en $O$, repère chaque avion par sa distance (en km) et par l'angle
qu'il fait avec la direction de l'est, celle de $\vec{u}$. L'avion $A$ est à $6$~km,
dans la direction $\frac{\pi}{2}$ ; l'avion $B$ est à $8$~km, dans la direction $0$.
\begin{enumerate}[label=\alph*)]
  \item Donner les affixes $z_{A}$ et $z_{B}$.
  \item Calculer la distance entre les deux avions.
\end{enumerate}

\partiefiche{Forme trigonométrique}

\exo[Vers la forme algébrique]
Écrire sous forme algébrique :
\begin{center}
$z_{1}=3\left(\cos\frac{\pi}{4}+i\sin\frac{\pi}{4}\right)$ \quad
$z_{2}=4\left(\cos\frac{\pi}{2}+i\sin\frac{\pi}{2}\right)$ \\
$z_{3}=7\left(\cos\pi+i\sin\pi\right)$ \quad
$z_{4}=6\left(\cos\left(-\frac{\pi}{6}\right)+i\sin\left(-\frac{\pi}{6}\right)\right)$ \\
$z_{5}=2\left(\cos\frac{2\pi}{3}+i\sin\frac{2\pi}{3}\right)$ \quad
$z_{6}=\sqrt{2}\left(\cos\frac{3\pi}{4}+i\sin\frac{3\pi}{4}\right)$
\end{center}

\exo[Lire, placer, calculer]
Soit $z_{1}=3\left(\cos\frac{5\pi}{6}+i\sin\frac{5\pi}{6}\right)$ et
$z_{2}=2\left(\cos\left(-\frac{\pi}{3}\right)+i\sin\left(-\frac{\pi}{3}\right)\right)$.
\begin{enumerate}[label=\alph*)]
  \item Donner le module et un argument de $z_{1}$ et de $z_{2}$.
  \item Placer les points d'affixes $z_{1}$ et $z_{2}$ (unité $1$~cm).
  \item Écrire $z_{1}$ et $z_{2}$ sous forme algébrique.
\end{enumerate}

\exo[Tensions de Fresnel]
En physique, une tension sinusoïdale d'amplitude $10$~V et de phase $\frac{\pi}{6}$ se
représente par $U_{1}=10\left(\cos\frac{\pi}{6}+i\sin\frac{\pi}{6}\right)$ ; une seconde
tension, de phase $-\frac{\pi}{6}$, par
$U_{2}=10\left(\cos\left(-\frac{\pi}{6}\right)+i\sin\left(-\frac{\pi}{6}\right)\right)$.
\begin{enumerate}[label=\alph*)]
  \item Écrire $U_{1}$ et $U_{2}$ sous forme algébrique.
  \item Vérifier que $U_{2}=\overline{U_{1}}$.
  \item Calculer $U_{1}+U_{2}$, puis l'amplitude $|U_{1}+U_{2}|$, arrondie au dixième.
\end{enumerate}

\partiefiche{De la forme algébrique à la forme trigonométrique}

\exo[Sur les axes]
Écrire sous forme trigonométrique : $z_{1}=6$, $z_{2}=4i$, $z_{3}=-2$ et $z_{4}=-3i$.

\exo[Vers la forme trigonométrique]
Écrire sous forme trigonométrique :
\begin{center}
$z_{1}=\sqrt{3}+i$ \quad $z_{2}=1-i\sqrt{3}$ \quad $z_{3}=-\sqrt{2}-i\sqrt{2}$ \\
$z_{4}=-3\sqrt{3}+3i$ \quad $z_{5}=2\sqrt{2}+2i\sqrt{2}$
\end{center}

\exo[Avec des $\sqrt{2}$]
\rappel{$\dfrac{1}{\sqrt{2}}=\dfrac{\sqrt{2}}{2}$ (on multiplie en haut et en bas par $\sqrt{2}$).}
Écrire sous forme trigonométrique : $w_{1}=1+i$, $w_{2}=3-3i$ et $w_{3}=-5+5i$.

\exo[Déphasage d'un dipôle]
Un dipôle a pour impédance $Z=30\sqrt{3}+30i$ (en ohms).
\begin{enumerate}[label=\alph*)]
  \item Donner la résistance $R=\mathrm{Re}(Z)$ et la réactance $X=\mathrm{Im}(Z)$.
  \item Calculer $|Z|$.
  \item L'argument $\varphi$ de $Z$ est le déphasage entre la tension et
    l'intensité. Calculer $\cos\varphi$ et $\sin\varphi$, puis donner $\varphi$.
  \item Écrire $Z$ sous forme trigonométrique.
  \item Sous une tension de $230$~V, l'intensité vaut $I=\dfrac{230}{|Z|}$. La calculer,
    arrondie au centième.
\end{enumerate}

\end{multicols}

\end{document}
